\chapter{光的衍射}

\problem{2.1}{一平行单色光照射到一小圆孔上，将其波面分成半波带，求第 $k$ 个带的半径。若极点到观察点的距离 $r_0$ 为 $1\,\mathrm{m}$，单色光波长为 $450\,\mathrm{nm}$，求此时第一半波带的半径。}
\answerarea

\problem{2.2}{一平行单色光从左向右垂直射到一个有圆形小孔的屏上，设此孔可以像摄像机光圈那样改变大小。问：（1）小孔半径应满足什么条件时，才能使得此小孔右侧轴线上距小孔中心 $4\,\mathrm{m}$ 的 $P$ 点的光强分别得到极大值和极小值；（2）$P$ 点最亮时，小孔直径应为多大？设此光的波长为 $500\,\mathrm{nm}$。}
\answerarea
\clearpage

\problem{2.3}{波长为 $500\,\mathrm{nm}$ 的单色点光源离光阑 $1\,\mathrm{m}$，光阑上有一个内外半径分别为 $0.5\,\mathrm{mm}$ 和 $1\,\mathrm{mm}$ 的透光圆环，接收点 $P$ 离光阑 $1\,\mathrm{m}$，求 $P$ 点的光强 $I$ 与没有光阑时的光强度 $I_0$ 之比。}
\answerarea

\problem{2.4}{波长为 $632.8\,\mathrm{nm}$ 的 $\mathrm{He\text{-}Ne}$ 激光射向直径为 $2.76\,\mathrm{mm}$ 的圆孔，与孔相距 $1\,\mathrm{m}$ 处放一屏。试问：（1）屏上正对圆孔中心的 $P$ 点是亮点还是暗点？（2）要使 $P$ 点变成与（1）相反的情况，至少要把屏幕分别向前或向后移动多少？}
\answerarea
\clearpage

\problem{2.5}{一波带片由五个半波带组成。第一半波带为半径 $r_1$ 的不透明圆盘，第二半波带是半径 $r_1$ 至 $r_2$ 的透明圆环，第三半波带是 $r_2$ 至 $r_3$ 的不透明圆环，第四半波带是 $r_3$ 至 $r_4$ 的透明圆环，第五半波带是 $r_4$ 至无穷大的不透明区域。已知 $r_1:r_2:r_3:r_4=1:\sqrt{2}:\sqrt{3}:\sqrt{4}$，用波长 $500\,\mathrm{nm}$ 的平行单色光照明，最亮的像点在距波带片 $1\,\mathrm{m}$ 的轴上。试求：（1）$r_1$；（2）像点的光强与入射光强之比；（3）光强极大值出现在轴上哪些位置上。}
\answerarea

\problem{2.6}{波长为 $\lambda$ 的点光源经波带片成一个像点，该波带片有 $100$ 个透明奇数半波带（$1,3,5,\cdots,199$），另外 $100$ 个不透明偶数半波带。比较用波带片和换上同样焦距和口径的透镜时该像点的强度比 $I:I_0$。}
\answerarea
\clearpage

\problem{2.7}{一平行单色光的波长为 $480\,\mathrm{nm}$，垂直照射到宽度为 $0.4\,\mathrm{mm}$ 的狭缝上，会聚透镜的焦距为 $60\,\mathrm{cm}$。分别计算当缝的两边到 $P$ 点的相位差为 $\pi/2$ 和 $\pi/6$ 时，$P$ 点离焦点的距离。}
\answerarea

\problem{2.8}{白光形成的单缝衍射图样中，其中某一波长的第三个次最大值与波长为 $600\,\mathrm{nm}$ 的光波的第二个次最大值重合，求该光波的波长。}
\answerarea
\clearpage

\problem{2.9}{波长为 $546.1\,\mathrm{nm}$ 的平行单色光垂直地射在 $1\,\mathrm{mm}$ 宽的缝上，若将焦距为 $100\,\mathrm{cm}$ 的透镜紧贴于缝的后面，并使光聚焦到屏上，试问衍射图样的中央到（1）第一最小值；（2）第一最大值；（3）第三最小值的距离分别为多少？}
\answerarea

\problem{2.10}{钠光通过宽 $0.2\,\mathrm{mm}$ 的狭缝后，投射到与缝相距 $300\,\mathrm{cm}$ 的照相底片上，所得的第一最小值与第二最小值间的距离为 $0.885\,\mathrm{cm}$，试问钠光的波长为多少？若改用 X 射线（$\lambda=0.1\,\mathrm{nm}$）做此实验，问底片上这两个最小值之间的距离是多少？}
\answerarea
\clearpage

\problem{2.11}{以纵坐标表示强度，横坐标表示屏上的位置，粗略地画出三缝的夫琅禾费衍射（包括缝与缝之间的干涉）图样。设缝宽为 $b$，相邻缝间的距离为 $d$，$d=3b$。注意缺级问题。}
\answerarea

\problem{2.12}{一束平行白光垂直入射在每毫米 $50$ 条刻痕的全息光栅上，问第一级光谱的末端和第二光谱的始端的衍射角 $\theta$ 之差为多少？（设可见光中最短的紫光波长为 $400\,\mathrm{nm}$，最长的红光波长为 $760\,\mathrm{nm}$。）}
\answerarea
\clearpage

\problem{2.13}{用可见光（$760\sim 400\,\mathrm{nm}$）照射全息光栅时，一级光谱和二级光谱是否重叠？二级和三级怎样？若重叠，则重叠范围是多少？}
\answerarea

\problem{2.14}{用波长为 $589\,\mathrm{nm}$ 的平行单色光照射一衍射光栅，其光谱的中央最大值和第二十级主最大值之间的衍射角为 $15^\circ 10'$，求该光栅 $1\,\mathrm{cm}$ 内的缝数。}
\answerarea
\clearpage

\problem{2.15}{用每毫米内有 $400$ 条刻痕的平面透射光栅观察波长为 $589\,\mathrm{nm}$ 的钠光谱，试问：（1）光垂直入射时，最多能观察到几级光谱？（2）光以 $30^\circ$ 角入射时，最多能观察到几级光谱？}
\answerarea

\problem{2.16}{白光垂直照射到一个每毫米 $250$ 条刻痕的平面透射光栅上，试问在衍射角为 $30^\circ$ 处会出现哪些波长的光？其颜色如何？}
\answerarea
\clearpage

\problem{2.17}{一波长为 $624\,\mathrm{nm}$ 的平行单色光照射一光栅，已知该光栅的缝宽 $b$ 为 $0.012\,\mathrm{mm}$，不透明部分的宽度 $a$ 为 $0.029\,\mathrm{mm}$，缝数 $N$ 为 $10^3$ 条。试求：（1）单缝衍射图样的中央角宽度；（2）单缝衍射图样中央宽度内能看到多少级光谱？（3）谱线的半宽度为多少？}
\answerarea

\problem{2.18}{$\mathrm{NaCl}$ 的晶体结构是简单的立方点阵，其相对分子质量 $M=58.5$，密度 $\rho=2.17\,\mathrm{g/cm^3}$，（1）试证相邻两离子间的平均距离为
\[\sqrt[3]{\frac{M}{2N_{\mathrm{A}}\rho}}=0.281\,9\,\mathrm{nm}\]
式中 $N_{\mathrm{A}}=6.02\times 10^{23}\,\mathrm{mol^{-1}}$ 为阿伏伽德罗常量；（2）用 X 射线照射晶面时，第二级光谱的最大值在掠射角为 $1^\circ$ 的方向上出现，试计算该 X 射线的波长。}
\answerarea
\clearpage

\problem{2.19}{波长为 $0.001\,47\,\mathrm{nm}$ 的平行 X 射线射在晶体界面上，晶体原子层的间距为 $0.28\,\mathrm{nm}$，试问光线与界面成什么角度时，能观察到二级光谱。}
\answerarea

\sproblemfig[0.34]{2.20}{如题 2.20 图所示有三条彼此平行的狭缝，宽度均为 $b$，缝距分别为 $d$ 和 $2d$，试用振幅矢量叠加法证明正入射时，夫琅禾费衍射强度公式为
\[I_\theta = I_0\frac{\sin^2 u}{u^2}\bigl[3+2(\cos 2v+\cos 4v+\cos 6v)\bigr]\]
式中 $u=\dfrac{\pi b\sin\theta}{\lambda}$，$v=\dfrac{\pi d\sin\theta}{\lambda}$。}{2_20}
\answerarea
\clearpage

\sproblemfig[0.24]{2.21}{一宽度为 $2\,\mathrm{cm}$ 的平面透射光栅上刻有 $12\,000$ 条刻痕。如题 2.21 图所示，以波长 $\lambda=500\,\mathrm{nm}$ 的单色光垂直投射，将折射率为 $1.5$ 的劈状玻璃片置于光栅前方，玻璃片的厚度从光栅的一端到另一端由 $1\,\mathrm{mm}$ 均匀变薄到 $0.5\,\mathrm{mm}$，试问第一级主最大方向改变了多少？}{2_21}
\answerarea

\sproblem{2.22}{一平行单色光投射于平面透射光栅上，其方向与光栅的法线成 $\theta_0$ 角，在和法线成 $11^\circ$ 和 $53^\circ$ 角的方向上出现第一级谱线，且位于法线的两侧。
\begin{enumerate}
  \item 试求入射角 $\theta_0$；
  \item 试问为什么在法线两侧能观察到一级谱线，而在法线同侧则能观察到二级谱线？
\end{enumerate}}
\answerarea
\clearpage

\sproblem{2.23}{一波长为 $600\,\mathrm{nm}$ 的平行单色光正入射到一平面透射光栅上，有两个相邻主最大分别出现在 $\sin\theta_1=0.2$ 和 $\sin\theta=0.3$ 处，第四级为缺级。
\begin{enumerate}
  \item 试求光栅常量；
  \item 试求光栅的缝可能的最小宽度；
  \item 在确定了光栅常量与缝宽之后，试列出在光屏上实际呈现的全部级数。
\end{enumerate}}
\answerarea

\sproblem{2.24}{研究性课题：设计实验测定 CD 或 DVD 光盘的光栅常量。}
\answerarea
\clearpage
